\documentclass[../../main.tex]{subfiles}
\begin{document}

\begin{flushright}
\footnotesize\textit{【自動文字起こし・要確認】}
\end{flushright}

[解]
(1) $[0, 1]$ では $0 < (1+x)^{-n}, x e^{x^2} \le e$ だから\\
（$x e^{x^2}$ は単調、$x=1$ とすれば $e$）
\[
0 \le (1+x)^{-n} \cdot x e^{x^2} \le e(1+x)^{-n}
\]
同じ区間で積分して
\[
0 \le b_n \le e \int_0^1 (1+x)^{-n} dx
\]
右辺を計算する。
\begin{align}
0 \le b_n &\le e \left[ \dfrac{1}{1-n}(1+x)^{-n+1} \right]_0^1 \\
&= \dfrac{e}{1-n} \left[ 2^{1-n} - 1 \right] \to 0 \quad (n \to \infty)
\end{align}
はさみうちから $b_n \to 0 \quad (n \to \infty)$

(2)
\begin{align}
a_n &= \left[ -\dfrac{1}{n}(1+x)^{-n} e^{x^2} \right]_0^1 + \dfrac{1}{n} \int_0^1 (1+x)^{-n} \cdot 2x \cdot e^{x^2} dx \\
&= -\dfrac{1}{n} \left[ e \cdot 2^{-n} - 1 \right] + \dfrac{2}{n} b_n
\end{align}
だから
\[
n a_n = \left( 1 - \dfrac{e}{2^n} \right) + 2 b_n \to 1 \quad (n \to \infty) \quad (\because (1))
\]
より
\[
n a_n \to 1 \quad (n \to \infty)
\]

\end{document}
